% ============================================================================
% Fichier  : partie0/P00_la_lettre.tex
% Titre    : Prologue : la Lettre d'Autorisation
% Auteur   : MINKA MI NGUIDJOÏ Thierry Emmanuel
% Version  : 1.0 -- novembre 2026
% ============================================================================
% Opération LOGONE : test d'intrusion de LiZbiZ par NYONG SÉCURITÉ, du 19 au
% 30 octobre 2026. Cas fictif, cohérent avec SEC3266 et GES4062. Adresses de
% documentation (RFC 5737) ou fictives.
% ============================================================================

\chapter{Prologue : la Lettre d'Autorisation}
\label{chap:prologue}
\chaptermark{Prologue}

\epigraph{%
    « Sans lettre d'autorisation, le meilleur test d'intrusion\\
    est une infraction. »%
}{Principe du cours}

\niveauChapitre{Fondamental}{Cours SEC3266 (Sécurité informatique).}

\begin{boxobjectifs}
\begin{itemize}
    \item Lire un mandat de test d'intrusion et en repérer les éléments
          essentiels.
    \item Situer le test dans la vie d'une organisation~: après une crise,
          après un audit, avant une décision.
    \item Distinguer ce que le test prouve de ce qu'il ne prouve pas.
    \item Séparer constatations, suppositions et opinions dans un résultat
          de test.
    \item Formuler les questions auxquelles le cours répond.
\end{itemize}
\end{boxobjectifs}

\bigskip

% ============================================================================
\section*{Question directrice : que faut-il signer avant de toucher un seul système~?}
% ============================================================================

Le 12~octobre 2026, M.~ETEME signe la lettre qui autorise le cabinet NYONG
SÉCURITÉ à attaquer le système d'information de LiZbiZ. Sept mois plus tôt,
un attaquant avait fait la même chose sans demander la permission
(Opération GAROUA). Cette fois, tout est écrit~: qui teste, quoi, quand,
comment, et qui appeler si quelque chose tourne mal. Ce prologue suit la
mission, appelée Opération LOGONE, de la signature au rapport.

\medskip
\textbf{Plan du chapitre.}
\begin{enumerate}
    \item Pourquoi un test, et pourquoi maintenant (section~\ref{sec:p00-pourquoi}).
    \item La lettre d'autorisation (section~\ref{sec:p00-lettre}).
    \item Douze jours de test (section~\ref{sec:p00-test}).
    \item Ce que le test a montré (section~\ref{sec:p00-resultats}).
    \item Les questions du cours (section~\ref{sec:p00-questions}).
\end{enumerate}

\separateur

% ============================================================================
\section{Pourquoi un Test, et Pourquoi Maintenant}
\label{sec:p00-pourquoi}
% ============================================================================

LiZbiZ\footnote{LiZbiZ, NYONG SÉCURITÉ, la BANQUE SANAGA, l'ANPP et leurs
personnels sont fictifs, comme les dates, les adresses et les montants du
cours~; les techniques, les normes et les textes de loi cités sont réels.}
a beaucoup changé depuis mars~: second facteur sur la passerelle, bastion,
sauvegardes immuables, registre des risques, audit d'octobre. Le comité des
risques veut savoir si ces mesures résistent à un attaquant réel, et non
seulement à une liste de vérification. Il commande un \termecle{test
d'intrusion}~: une attaque autorisée, bornée et restituée.

\begin{tableaucours}{L{3.4cm}L{8.8cm}}{%
    \tableauHeader{Acteur} &
    \tableauHeader{Rôle dans l'Opération LOGONE}
}
M.~ETEME & Directeur général de LiZbiZ, signataire de la lettre
d'autorisation \\
M.~KOUNOU & Responsable informatique, contact technique d'urgence \\
Mme~MOUAFO & Gestionnaire des risques, commanditaire opérationnelle \\
M.~BIKORO & Chef de mission, NYONG SÉCURITÉ \\
Mme~NGO~MBOCK & Pentesteuse, tests internes et Active Directory \\
M.~TAMO & Analyste, tests web et ingénierie sociale \\
\end{tableaucours}
\vspace{-0.8em}
\begin{center}{\small\itshape Tableau~0.1. Les acteurs de l'Opération LOGONE}\end{center}

% ============================================================================
\section{La Lettre d'Autorisation}
\label{sec:p00-lettre}
% ============================================================================

\begin{tableaucours}{L{3.4cm}L{8.8cm}}{%
    \tableauHeader{Rubrique} &
    \tableauHeader{Contenu de la lettre du 12 octobre}
}
Parties & LiZbiZ SA (M.~ETEME) et NYONG SÉCURITÉ (M.~BIKORO) \\
Objectif & Vérifier qu'un attaquant externe ou interne ne peut pas obtenir
les droits d'administrateur du domaine ni lire la base des clients \\
Périmètre externe & Plage 198.51.100.0/27 et domaine lizbiz.cm \\
Périmètre interne & Réseaux de Douala (10.10.0.0/16) et de Garoua
(10.20.0.0/16), sauf 10.10.50.0/24 (paie, hébergée pour un tiers) \\
Exclusions & Déni de service~; système de paie~; intrusion physique \\
Ingénierie sociale & Courriels d'hameçonnage simulés à 60~adresses validées
par la direction des ressources humaines \\
Fenêtre & Du 19 au 30~octobre~; tests internes de 8~h à 18~h, sauf accord
écrit \\
Contacts d'urgence & M.~KOUNOU et M.~BIKORO, joignables jour et nuit~; mot
d'arrêt convenu \\
Données & Aucune copie de données personnelles~; preuves masquées~;
destruction à la remise du rapport \\
\end{tableaucours}
\vspace{-0.8em}
\begin{center}{\small\itshape Tableau~0.2. Extrait de la lettre d'autorisation de l'Opération LOGONE}\end{center}

\begin{boxalert}
\textbf{Violation typique.} Tester une adresse « voisine » du périmètre
parce qu'elle semble appartenir au client.

\smallskip
\textbf{Règle~:} seul ce qui est écrit est autorisé. Le réseau de la paie,
hébergé pour un tiers, n'appartient pas au signataire~: le tester, même par
erreur, serait un accès non autorisé (Chapitre~\ref{chap:droit}).
\end{boxalert}

% ============================================================================
\section{Douze Jours de Test}
\label{sec:p00-test}
% ============================================================================

\begin{tableaucours}{L{2.4cm}L{9.8cm}}{%
    \tableauHeader{Dates} &
    \tableauHeader{Activité}
}
19--20~oct. & Reconnaissance externe~: sources ouvertes, sous-domaines,
services exposés \\
21--22~oct. & Tests externes~: passerelle, portail de suivi, messagerie \\
23~oct. & Campagne d'hameçonnage simulée \\
26--28~oct. & Tests internes depuis un poste fourni~: balayage, Active
Directory, partages \\
29~oct. & Démonstration du chemin vers les administrateurs du domaine~;
arrêt des tests \\
30~oct. & Réunion de restitution à chaud \\
\end{tableaucours}
\vspace{-0.8em}
\begin{center}{\small\itshape Tableau~0.3. Le calendrier de l'Opération LOGONE}\end{center}

% ============================================================================
\section{Ce que le Test a Montré}
\label{sec:p00-resultats}
% ============================================================================

\begin{boxscenario}{Opération LOGONE. La restitution du 30 octobre}
\textbf{Constatations.}
\begin{itemize}
    \item Externe~: la passerelle exige bien un second facteur~; aucun accès
          obtenu depuis Internet.
    \item Hameçonnage~: 14~clics sur 60~destinataires, 5~saisies de mot de
          passe, 9~signalements au support.
    \item Interne~: en 5~heures, depuis un poste ordinaire, Mme~NGO~MBOCK
          obtient les droits d'administrateur du domaine en exploitant le
          mot de passe faible d'un compte de service.
    \item Détection~: le centre de surveillance de nuit signale le balayage
          interne en 40~minutes, mais ne voit pas l'obtention des droits.
    \item Bilan~: 14~constats, dont 2 critiques, 4 élevés, 5 moyens et
          3 faibles.
\end{itemize}

\medskip
\textbf{Analyse} (confiance élevée). La porte d'entrée de mars est fermée~;
la maison, elle, reste ouverte à qui y entre. Le test ne prouve pas
qu'aucune autre voie n'existe~: il prouve que celle-ci existait, à cette
date, dans ce périmètre.
\end{boxscenario}

\begin{boxmetier}{Pentesteur}
\textbf{Arrêter quand l'objectif est démontré.} Mme~NGO~MBOCK s'arrête
dès qu'elle obtient les droits~: elle capture la preuve (une capture
d'écran du groupe, sans données), prévient M.~BIKORO, qui prévient
M.~KOUNOU. Aller plus loin n'apporterait rien et créerait un risque.
\end{boxmetier}

\begin{boxmetier}{RSSI}
\textbf{Commander un test pour décider.} Le RSSI commande un test avec une
question précise (« peut-on devenir administrateur du domaine~? »), et non
pour obtenir une note. Le rapport doit l'aider à décider où investir.
\end{boxmetier}

% ============================================================================
\section{Les Questions du Cours}
\label{sec:p00-questions}
% ============================================================================

\begin{tableaucours}{L{3.4cm}L{6.6cm}L{2.2cm}}{%
    \tableauHeader{Partie} &
    \tableauHeader{Question posée par l'Opération LOGONE} &
    \tableauHeader{Chapitres}
}
I. Cadre & Qu'est-ce qui rend ce test légal et utile~? & \ref{chap:ethique}
à~\ref{chap:methodes} \\
II. Reconnaître & Qu'a-t-on appris sans toucher aux systèmes~? &
\ref{chap:osint} à~\ref{chap:vulnerabilites} \\
III. Systèmes & Comment un compte de service mène-t-il au domaine~? &
\ref{chap:exploitation} à~\ref{chap:mdp} \\
IV. Applications & Le portail et l'API sont-ils sûrs~? & \ref{chap:web}
à~\ref{chap:nuage} \\
V. Au-delà & Les personnes et le SOC tiennent-ils~? & \ref{chap:sansfil}
à~\ref{chap:detection} \\
VI. Restituer & Comment écrire, corriger, vérifier, mesurer~? &
\ref{chap:rapport} à~\ref{chap:mesurer} \\
\end{tableaucours}
\vspace{-0.8em}
\begin{center}{\small\itshape Tableau~0.4. Les questions du cours}\end{center}

\separateur

\begin{boxresume}
\textbf{Ce qu'il faut retenir de ce chapitre~:}
\begin{itemize}
    \item \textbf{Un test d'intrusion} est une attaque autorisée, bornée et
          restituée.
    \item \textbf{La lettre d'autorisation} fixe parties, objectif,
          périmètre, exclusions, fenêtre, contacts et traitement des données.
    \item \textbf{Seul ce qui est écrit est autorisé}~; un tiers non
          signataire est hors périmètre.
    \item \textbf{On s'arrête} quand l'objectif est démontré.
    \item \textbf{Un test prouve} l'existence d'une voie, pas l'absence
          d'autres voies.
\end{itemize}
\end{boxresume}

\begin{boxinfo}
\textbf{Points de contrôle de l'élève.} Je sais~: énumérer les rubriques
d'une lettre d'autorisation~; repérer une cible hors périmètre~; résumer
les résultats de l'Opération LOGONE~; dire ce qu'un test prouve et ne
prouve pas.
\end{boxinfo}

\bigskip

\begin{boxexo}[title={\ding{45}\quad Exercice~0.1. Le taux \niveaubase}]
Calculez le taux de clics, de saisies et de signalements de la campagne
d'hameçonnage.
\end{boxexo}

\begin{boxexo}[title={\ding{45}\quad Exercice~0.2. Hors périmètre \niveaubase}]
Les adresses 10.10.50.12, 10.20.3.50 et 198.51.100.40 sont-elles dans le
périmètre de la lettre~? Justifiez.
\end{boxexo}

\begin{boxexo}[title={\ding{45}\quad Exercice~0.3. Les registres \niveaubase}]
Classez en constatation, supposition ou opinion~: « le SOC est
incompétent »~; « le balayage a été signalé en 40~minutes »~; « un
attaquant réel aurait été plus discret ».
\end{boxexo}

\begin{boxexo}[title={\ding{45}\quad Exercice~0.4. Prouver \niveaumoyen}]
Pourquoi le test ne prouve-t-il pas que LiZbiZ est sûre depuis Internet~?
Que peut-on en dire honnêtement~?
\end{boxexo}

\begin{boxexo}[title={\ding{45}\quad Exercice~0.5. L'arrêt \niveaumoyen}]
Mme~NGO~MBOCK aurait pu lire la base des clients après avoir obtenu les
droits. Pourquoi s'arrête-t-elle~? Quelle preuve garde-t-elle~?
\end{boxexo}

\begin{boxexo}[title={\ding{45}\quad Exercice~0.6. Le tiers \niveaumoyen}]
Le système de paie est hébergé par LiZbiZ pour un tiers. Que faudrait-il
pour l'inclure dans un prochain test~?
\end{boxexo}

\begin{boxexo}[title={\ding{45}\quad Exercice~0.7. Les données \niveaumoyen}]
Pourquoi la lettre interdit-elle de copier des données personnelles~? Quels
risques courrait le cabinet sans cette clause~?
\end{boxexo}

\begin{boxexo}[title={\ding{45}\quad Exercice~0.8. La détection \niveaumoyen}]
Le SOC a vu le balayage mais pas l'obtention des droits. Qu'en concluez-vous
pour ses règles de détection~?
\end{boxexo}

\begin{boxexo}[title={\ding{45}\quad Exercice~0.9. Une lettre \niveauexpert}]
Rédigez la lettre d'autorisation complète de l'Opération LOGONE, à partir
du Tableau~0.2.
\end{boxexo}

\begin{boxexo}[title={\ding{45}\quad Exercice~0.10. La question \niveauexpert}]
Proposez trois questions précises que la BANQUE SANAGA pourrait poser à un
test d'intrusion, et pour chacune le périmètre nécessaire.
\end{boxexo}

\bigskip

\begin{boxinfo}
\textbf{Pour aller plus loin.}
\begin{itemize}
    \item NIST SP~800-115~\cite{nist800115}.
    \item Loi n°~2010/012, articles~65 et suivants~\cite{loi2010012}.
\end{itemize}
\end{boxinfo}

\begin{boxinfo}
\textbf{Transition.} Avant les techniques, les mots et les métiers. Le
Chapitre~\ref{chap:ethique} définit le hacking éthique et les types de
tests.
\end{boxinfo}
